\documentclass[12pt]{article}
\usepackage[utf8]{inputenc}
\usepackage[T1]{fontenc}
\usepackage{hyperref}
\title{Kembra Language Resources for K-12 Curriculum Development}
\author{Compiled by Claude Sonnet 5.5 (claude-sonnet-5-5) for Andr\'as Kornai}
\date{}
\begin{document}
\maketitle

\section{Introduction}
Kembra (ISO 639-3 xkw) is a South Pauwasi language of Western New Guinea, now placed in the Lepki-Murkim-Kembra grouping. It is seriously endangered and is no longer passed on to children. Sources disagree on where it is spoken. Wikipedia gives Kiambra village (Keerom Regency) and about twenty speakers. Glosbe and Ethnologue-derived pages give Jayawijaya Regency. In my searches I found no school materials, orthography, primer or published dictionary.

\section{Primary Educational Resources}
None found. I found no textbook, primer or school programme for Kembra.

\section{Linguistic and Reference Materials}
\begin{itemize}
\item Wikipedia, ``Kembra language'': an overview of classification, location, tone and documentation status. It says the data remain unpublished in the field notes of Barnabas Konel and Roger Doriot. \url{https://en.wikipedia.org/wiki/Kembra_language}
\item ASJP (Automated Similarity Judgment Program) entry for Kembra, which holds a basic word list. \url{https://asjp.clld.org/languages/KEMBRA}
\item Glosbe language page for xkw. I did not check its content, so its value is unverified. \url{https://glosbe.com/lib/xkw}
\end{itemize}

\section{Religious and Community Resources}
None found in the search results.

\section{Digital and Audio Resources}
\begin{itemize}
\item Endangered Languages Project entry for Kembra, with a status assessment and metadata. \url{https://endangeredlanguages.com/lang/1897?hl=en}
\end{itemize}
No audio or video recordings were found.

\section{Curriculum Development Context}
Very little is documented, and the language is moribund. A school curriculum in Kembra is not feasible without prior documentation work with the remaining speakers. The first steps would be to record elders, settle on an orthography (the language is tonal) and publish the existing field data with the community's consent. Any curriculum would most realistically be a language-revitalisation programme in Indonesian.
\end{document}